%==============================================================================
\chapter{Analysis \& correction design}\label{ch:analysis}
%==============================================================================

% TODO(screenshot): the Analysis view: loaded measurement set, the transfer-function
% plot (thin individual curves, thick average, correction, corrected response),
% magnitude and phase; the smoothing, correction-level, boost-ceiling and
% analysis-range controls; the subwoofer-integration controls; the export buttons.
\screenshot{analysis-view.png}{The Analysis view: individual smoothed transfer
functions (thin), their average (thick), the proposed correction and the
corrected response, in magnitude and phase; with the nth-octave smoothing,
correction-level, boost-ceiling and analysis-range controls, the subwoofer
integration section (crossover, Invert, Mains delay), and the Export buttons.}

The Analysis view turns the measurements of one loudspeaker into a correction filter. You load the measurement set, check the averaged response on the plot, tune the \gterm{smoothing}{smoothing} and the correction controls until the corrected curve looks right, optionally align a subwoofer, then export the filter and assign it to an output. This chapter describes what each step does and how to set it. The signal processing behind it, with the equations, is in Appendix~\ref{ch:anatheory}.

\section{Loading and estimating the response}\label{sec:welch}
Click "Load measurements\ldots" and select the files of one loudspeaker, measured at several microphone positions around the listening point. For each file SuperMoTo estimates the loudspeaker-room \gterm{tf}{transfer function} from the stimulus and the microphone signal stored in it (Section~\ref{sec:export}). If a microphone calibration is in effect, it is divided out at this stage (Section~\ref{sec:micsource}), so the curves describe the loudspeaker and the room rather than the microphone.

\subsection{Choosing the estimator}
The "TF" selector chooses how each transfer function is estimated.

\begin{itemize}[nosep]
  \item \textbf{\gterm{welch}{Welch}} averages the spectra of many overlapping segments of the recording. It works with any stimulus and is always available.
  \item \textbf{Sweep (Farina)} deconvolves a logarithmic sweep. It is offered, and selected automatically, as soon as the loaded folder's \texttt{measurement.xml} records the sweep's parameters (Section~\ref{sec:manifest}). It recovers the whole band with its true phase from a single sweep, and it moves the harmonic distortion out of the linear response instead of leaving it folded in, which gives a cleaner estimate.
\end{itemize}

\noindent Everything after the estimate, from averaging to export, is identical for both.

\subsection{The Welch window}
The "Welch window" sets the length of the analysis segments, from 16384 to 262144 samples, 65536 by default. A longer window gives a finer frequency resolution, about \SI{0.67}{\hertz} at \SI{44.1}{\kilo\hertz} with the default, which is what resolves low-frequency \gterm{roommode}{room modes}. A shorter window cuts the same recording into more segments to average, which lowers the noise of the estimate. With the swept-sine estimator the window also sets how much of the deconvolved response is kept, and a window too large for the sweep makes the analysis fall back to Welch (Section~\ref{sec:sweep}).

\subsection{Delay removal and spatial averaging}\label{sec:spatialavg}
Each microphone position is at a different distance from the loudspeaker, so each measurement arrives with a different delay. SuperMoTo finds that delay from the peak of each measurement's \gterm{ir}{impulse response} and removes it, so every measurement starts at the same instant. The responses are then averaged as complex quantities, magnitude and phase together. This \gterm{spatialavg}{spatial averaging} is what makes several positions worthwhile. The \gterm{comb}{comb filtering} that reflections produce, narrow peaks and nulls that move with the microphone, tends to cancel in the average, leaving the response the loudspeaker reproduces consistently across the listening area~\cite{Toole,Mourjopoulos1994}. A correction designed from a single position would instead try to fill nulls that do not exist a few centimetres away. Appendix~\ref{ch:spatial} discusses how many positions to take and how far apart.

\subsection{The swept-sine estimator}\label{sec:sweep}
The sweep method needs the sweep's parameters, which SuperMoTo reads from the measurement folder's manifest. Three practical consequences follow.

\begin{itemize}[nosep]
  \item The sweep rate is quantized, so the actual sweep is slightly longer or shorter than the "Duration" you asked for. It sounds the same, and the quantization is what lets the harmonic orders come out with their true phase (Section~\ref{sec:sweepharm}).
  \item Harmonic distortion orders 2 to 5 appear on the plot as the H2\ldots{}H5 traces. They follow the same \gterm{smoothing}{smoothing} and microphone calibration as the other curves, and are averaged across positions in power, since distortion phase is not coherent from one position to the next.
  \item An order is dropped when the sweep is too short to separate it cleanly, which is why a short measurement shows fewer distortion curves than a long one.
\end{itemize}

\noindent A file is estimated with Welch instead whenever

\begin{itemize}[nosep]
  \item the folder carries no sweep identity, as with files loaded by hand, a noise run, or a folder recorded before the manifest stored it;
  \item the recording is shorter than the sweep it claims to be;
  \item the Welch window is too large for the sweep, so the harmonic responses would not stay clear of the linear one.
\end{itemize}

\noindent The fallback is per file and silent. A set that mixes sweep and noise runs still analyses, each file by whichever estimator applies to it. Welch remains the right, and only, choice for a noise stimulus, and stays selectable for a sweep recording if you want to compare the two.

\section{What the plot shows}
\begin{itemize}[nosep]
  \item the \textbf{individual} smoothed \gterm{tf}{transfer functions}, one per position (thin lines);
  \item their complex \textbf{average} (thick line);
  \item the proposed \textbf{correction};
  \item the \textbf{corrected response}, the average with the correction applied;
  \item the \textbf{subwoofer} response, when a subwoofer set is loaded (Section~\ref{sec:subalign});
  \item the \textbf{harmonic-distortion} magnitudes H2\ldots{}H5, with the swept-sine estimator only.
\end{itemize}

\noindent By default, magnitudes are normalised to the \SIrange{200}{2000}{\hertz} mean of the average, so \SI{0}{\dB} is the mid-band level, independent of the absolute measurement gain. Phase is shown as principal value with the propagation delay already removed. The correction's phase follows the "Phase" setting, so with "Min phase" selected it shows the phase the exported filter will actually have. Two selectors change what the panel displays.

\subsection{Level reference}\label{sec:levelref}
The "Level" selector sets the vertical reference of the measured curves. It is a display choice only. The correction is designed from the same data either way, and the correction curve itself is always shown in absolute \si{\dB}, since it is a gain and has no reference to choose.

\begin{itemize}[nosep]
  \item \textbf{Normalized} is the default described above, with \SI{0}{\dB} at the \SIrange{200}{2000}{\hertz} mean. Best for judging shape, and the only sensible choice when comparing speakers measured at different gains.
  \item \textbf{Absolute dB} shows the recorded level per unit of stimulus level, with no normalisation. Use it to compare the actual sensitivity of two speakers, or to see how much level a correction is giving away.
  \item \textbf{dB SPL} shows the estimated sound pressure during the measurement. This needs both an SPL calibration and the run's stimulus level, so it is offered only when the folder's manifest supplies them (Section~\ref{sec:manifest}), and greys out otherwise. It is exact for a sweep run and approximate for noise, the stimulus level being defined for a sine.
\end{itemize}

\subsection{Frequency response or impulse response}\label{sec:irview}
The "View" selector swaps the frequency plot for a time-domain view of the same speaker, at the current "\gterm{fir}{FIR} length":

\begin{itemize}[nosep]
  \item "measured" is the averaged response rendered back to the time domain;
  \item "corrected" is what the loudspeaker is predicted to do once the exported filter is loaded, following the "Phase" setting;
  \item "+ sub (as measured)", with a subwoofer loaded, is the corrected main summed with the subwoofer at the level it was measured, with its \gterm{polarity}{polarity} and the "Mains delay" applied.
\end{itemize}

\noindent Time zero is placed at the centre of the buffer, where a \gterm{linphase}{linear-phase} correction puts its main peak. This is the quickest way to see what the length and phase settings are really doing. It shows how much the corrected response rings, whether the FIR is long enough for the correction to decay inside it rather than being truncated, and how much pre-ringing a linear-phase filter introduces. Changing "FIR length" or "Phase" updates it immediately.

\section{Fractional-octave smoothing}\label{sec:smoothing}
\gterm{smoothing}{Smoothing} is applied both to the displayed curves and to the average the correction is derived from. It averages the complex response, magnitude and phase together, over a band a fixed fraction of an octave wide around each frequency, so it is narrow in the bass and wide in the treble. That matches the ear's roughly logarithmic frequency resolution and keeps the correction from chasing the dense, narrow interference structure of the room~\cite{Hatziantoniou2000}. The available fractions are Off, 1/24, 1/12, 1/6, 1/3, 1/2 and 1~octave.

\subsection{Frequency-dependent smoothing}\label{sec:varsmooth}
The "Smooth LF/HF" controls set two fractions independently, 1/6~octave each by default. The low-frequency one applies at and below \SI{100}{\hertz}, the high-frequency one at and above \SI{10}{\kilo\hertz}, and the fraction is blended smoothly in between (Section~\ref{sec:thsmooth}). Setting both to the same value gives uniform smoothing. This lets you correct the bass in detail and the treble only broadly, in a single correction.

\subsection{Why smoothing matters, and a recommendation}\label{sec:smoothrec}
\textbf{Do not try to correct a raw or lightly-smoothed response.} An unsmoothed in-room measurement is full of sharp, deep nulls and peaks produced by reflections and standing waves. These narrow features have three properties that make inverting them counter-productive:
\begin{itemize}[nosep]
  \item They are largely non-\gterm{minphase}{minimum-phase}. These are true cancellations that no magnitude or phase filter can undo, and attempting it only demands enormous boosts that the boost ceiling (Section~\ref{sec:correction}) then has to claw back, wasting headroom~\cite{NeelyAllen,Mourjopoulos1994}.
  \item They are position-dependent. A null at the measurement point has moved or vanished a few centimetres away, so "filling" it corrects one point while degrading every other listening position.
  \item Inverting a narrow notch produces a long, high-$Q$ filter whose ringing (pre-/post-echo in the rendered \gterm{fir}{FIR}) is itself audible.
\end{itemize}
The correction here is derived from the smoothed complex average, so the smoothing setting is your primary control over how much of this fine structure the inverse is even allowed to see. Combined with the \gterm{spatialavg}{spatial averaging} of Section~\ref{sec:spatialavg}, which already cancels much of the position-dependent \gterm{comb}{comb filtering}, enough smoothing keeps the correction on the broad, reproducible trends of the loudspeaker-room response, exactly the part that is stable across the listening area and safe to invert.

\paragraph{Practical guidance.}
\begin{itemize}[nosep]
  \item Treat the corrected curve, not the raw one, as the target. It should look smooth, not a mirror image of every wiggle in the measurement. If the correction curve is spiky, increase the smoothing.
  \item Exploit the frequency-dependent smoothing (Section~\ref{sec:varsmooth}). It is generally safe to correct more aggressively at low frequencies, which are broad, modal and relatively position-stable, and only lightly at high frequencies. There the response is dominated by direct sound plus early reflections, so only broad trends should be corrected and narrow dips left alone~\cite{Toole}. A good starting point is a finer LF fraction (e.g.\ 1/12-1/6~octave) with a coarser HF fraction (e.g.\ 1/3-1~octave).
  \item For uniform smoothing, set both controls equal, with 1/6~octave on both as a sensible neutral baseline. If a correction sounds harsh, "phasey" or over-processed, move coarser and/or lower the "Correction level" (Section~\ref{sec:correction}) rather than fighting individual peaks.
  \item Use the "Range" to keep the correction off the extreme top and bottom entirely, where even heavy smoothing cannot make the data trustworthy.
  \item Smoothing is your first line of \gterm{regularization}{regularisation}. The boost ceiling and the fade outside the analysis range are the safety nets, not a substitute for choosing a sensible smoothing.
\end{itemize}

\section{Designing the correction}\label{sec:correction}
The correction compensates both magnitude and phase. It starts from the exact inverse of the smoothed average, normalised to its \SIrange{200}{2000}{\hertz} mean level, which would flatten the response completely. Exact inversion is only safe for the \gterm{minphase}{minimum-phase} part of the response, because deep room nulls cannot be filled without enormous, ringing boosts~\cite{NeelyAllen,Mourjopoulos1994}. Four stages \gterm{regularization}{regularise} that inverse into the correction that is actually applied, and Section~\ref{sec:thcorrection} gives their formulas.

\begin{itemize}[nosep]
  \item \textbf{Max boost} limits how much the correction may lift the response, from 0 to \SI{24}{\dB} with a default of \SI{12}{\dB}. The limit is approached through a soft knee rather than a hard clamp, so notches do not leave a kink. Cuts are never limited, because attenuating a peak is always safe and only filling a null is dangerous.
  \item \textbf{Low-frequency slope.} The correction rolls off below \SI{30}{\hertz} like a second-order high-pass, so it does not chase the woofer's natural roll-off by boosting subsonic content where the speaker produces only noise. This stage has no control.
  \item \textbf{Correction level} morphs from no correction at 0 to the full correction at 1, the default, scaling magnitude and phase together. Intermediate values are useful when full inversion sounds over-processed.
  \item \textbf{Range} sets the band the correction acts on, \SI{20}{\hertz} to \SI{20}{\kilo\hertz} by default, with editable presets. Outside it the correction fades back to unity over half an octave, so noise and content the loudspeaker cannot reproduce are neither inverted nor written into the exported filter.
\end{itemize}

\begin{figure}[ht]
\centering
\begin{tikzpicture}
\begin{axis}[
  width=0.85\textwidth, height=6cm,
  xmode=log, xlabel={Frequency (Hz)}, ylabel={Magnitude (dB)},
  xmin=20, xmax=20000, ymin=-18, ymax=18,
  grid=both, legend pos=south west, legend cell align=left,
  title={Illustrative correction principle (schematic, not measured data)}]
% Measured average: a bump + a notch (schematic)
\addplot[thick, smtblue, smooth] coordinates {
  (20,-8)(40,-3)(70,2)(120,6)(200,3)(400,0)(800,-1)(1500,4)
  (3000,-6)(6000,1)(12000,-2)(20000,-5)};
\addlegendentry{measured average}
% Correction: inverse, with LF slope + capped boost
\addplot[thick, smtgold, smooth, densely dashed] coordinates {
  (20,6)(40,4)(70,-2)(120,-6)(200,-3)(400,0)(800,1)(1500,-4)
  (3000,6)(6000,-1)(12000,2)(20000,3)};
\addlegendentry{correction (inverse, capped)}
% Corrected = roughly flat
\addplot[thick, smtgreen, smooth] coordinates {
  (20,-2)(40,0)(70,0)(120,0)(200,0)(400,0)(800,0)(1500,0)
  (3000,0)(6000,0)(12000,0)(20000,-2)};
\addlegendentry{corrected response}
\end{axis}
\end{tikzpicture}
\caption{Schematic of the correction principle: the correction (gold, dashed)
is the regularised inverse of the smoothed measured average (blue); applying
it yields an approximately flat corrected response (green). The boost ceiling
prevents the correction from fully inverting deep notches; the low-frequency
slope rolls the correction off below the target corner. \emph{This plot is
illustrative, not real measurement data.}}
\label{fig:correction}
\end{figure}

\section{Subwoofer phase alignment}\label{sec:subalign}
Where a main and a subwoofer overlap around the \gterm{crossover}{crossover}, their outputs add as complex pressures. The sum is only flat if the two are in phase through the crossover region, otherwise it shows a dip, or a deep cancellation if they are in anti-phase~\cite{Linkwitz1976}. SuperMoTo can steer the corrected main's phase onto the subwoofer's so the pair sums coherently, without changing either magnitude.

\subsection{Setting it up}
\begin{enumerate}[nosep]
  \item Load the main's measurement set as usual.
  \item Click "Load sub measurements\ldots" and select the subwoofer set, measured at the same positions and loaded in the same order as the main's. Each subwoofer measurement is paired with the main measurement at the same position, which preserves their relative timing (Section~\ref{sec:thsubalign}).
  \item Set the "Crossover" to the crossover frequency of the system, \SI{80}{\hertz} by default. Below it the main's phase is steered onto the subwoofer's, and the steering is released over the octave above, so the main keeps its own flat-phase correction in its passband.
  \item Tick "Invert sub" if the subwoofer is wired, or behaves, out of \gterm{polarity}{polarity}.
  \item Tune the "Mains delay", described below, and export.
\end{enumerate}

\subsection{The Mains delay control}\label{sec:timealign}
The subwoofer usually lags the mains by several milliseconds, its own \gterm{groupdelay}{group delay} added to any difference in path or processing latency. Folding all of that into the phase correction would demand a long, low-frequency-heavy \gterm{fir}{FIR}. Instead the "Mains delay" slider, \SI{-40}{} to \SI{40}{\milli\second}, declares a bulk delay that you apply to the main output(s) with the per-output delays in the matrix, and the correction only carries the residual. A positive value delays the main(s) and a negative one delays the subwoofer. When it matches the real offset, the corrected main's phase is already flat through the crossover, which gives a short, well-behaved FIR. The corrected-phase curve visibly flattens as you approach the right value.

On loading the subwoofer set, SuperMoTo estimates the offset from the slope of the subwoofer's phase around the crossover and shows it as a recommendation. It is not applied automatically, because the value is a promise that the matching delay will be set in the matrix. You opt in by moving the slider, and the on-screen message then states the exact delay to dial in, and on which output. Appendix~\ref{ch:delays} explains why the phase slope, rather than the arrival time, is the right measure for a subwoofer.

\paragraph{Applying it with the FIR.} To skip the manual matrix step, tick "Apply bulk delay". "Export correction IR" then writes the delay onto the assigned output at the same time it loads the FIR, so the bulk delay and the residual correction land together. Output delays cannot be negative, so a negative value is clamped to zero, and in that case you delay the subwoofer's output by hand. This is the per-speaker measured route of Section~\ref{sec:ambisonics}. Each main's exported delay equals its path-length term, so all mains, and the subwoofer, arrive together at the centre.

The alignment only exists in the \gterm{linphase}{linear-phase} rendering. With "Min phase" selected the exported filter carries no phase steering, and "Crossover", "Invert sub" and "Mains delay" have no effect on it (Section~\ref{sec:render}).

\section{Phase type and latency}\label{sec:render}
Both exports turn a complex spectrum into a finite \gterm{ir}{impulse response} of the chosen "\gterm{fir}{FIR} length", from 256 to 65536 samples, 4096 by default. Longer filters reach lower frequencies, at the cost of more latency in linear phase. The "Phase" selector chooses how the correction is rendered, while the measured-response export is always linear phase.

\begin{itemize}[nosep]
  \item \textbf{Linear phase}, the default, keeps the full correction, magnitude and phase, including the subwoofer alignment. Its energy is centred in the filter, which adds a latency of half the FIR length on that output. The matrix's \gterm{latencycomp}{inter-output latency compensation} absorbs it automatically so all outputs stay time-aligned (Chapter~\ref{ch:matrix}).
  \item \textbf{Min phase} rebuilds the filter from the correction's magnitude alone, as a \gterm{minphase}{minimum-phase} filter (Section~\ref{sec:minphase}). The magnitude correction is identical, and the filter adds essentially no latency, which suits monitoring while tracking. It gives up the \gterm{allpass}{excess-phase} flattening and the subwoofer alignment, since a minimum-phase filter cannot impose an arbitrary phase.
\end{itemize}

\noindent The latency compensation reads each FIR's actual peak position, so it recognises both renderings without any setting. It aligns every fed output to the longest engaged FIR, however, so the low-latency benefit only appears when all engaged outputs use minimum-phase FIRs or none. A single \gterm{linphase}{linear-phase} FIR in the active configuration pulls the rest back up to its delay.

\section{Exporting}
\begin{itemize}[nosep]
  \item \textbf{Export IR} saves the measured response itself, the delay-aligned complex average, band-limited to the analysis range so out-of-band room and microphone noise does not leak into the file.
  \item \textbf{Export correction IR} renders the correction to a mono 32-bit float wav with the chosen "\gterm{fir}{FIR} length" and "Phase", and can load it directly onto the output selected in "Assign to". In linear phase it adds half the FIR length of delay on that output, which the \gterm{latencycomp}{inter-output latency compensation} absorbs. In \gterm{minphase}{minimum phase} it adds essentially none.
\end{itemize}
